\begin{dcomment}{context}{p1:Le fait que $b$ soit un plus grand élément}{What an ordinary pushout forgets}
Grothendieck is looking for categorical gluings that preserve a homotopy type. The little span that appears throughout this appendix is a particularly severe test. Its indexing category
\[
  \Psi=(b_0\longleftarrow a_0\longrightarrow c_0)
\]
has an initial object, hence a contractible nerve. Nevertheless, taking the ordinary colimit of its representable presheaves can introduce higher homotopy into the category of elements. The purpose of the examples is to locate this discrepancy precisely enough to decide which inclusions can be used in a homotopy theory on \(\Cat\).

We write \(BC=|NC|\), and interpret \(\Winf\) as the functors inducing weak homotopy equivalences on these classifying spaces. In this terminology ``aspherical'' means weakly contractible, not merely a space with vanishing higher homotopy groups. For a presheaf \(F:C^{\mathrm{op}}\to\Ens\), the notation \(C/F\) denotes its category of elements: its objects are \((c,x)\), and an arrow \((c,x)\to(c',x')\) is an arrow \(u:c\to c'\) with \(F(u)(x')=x\). Thus a represented presheaf has contractible category of elements \(C/c\). The central phenomenon is that a colimit of represented presheaves need not retain that property.

This is already the tension described in \citet[\S69]{Grothendieck1983}: categorical integration gives homotopy colimits, whereas actual amalgamated sums require additional care. The letter to Thomason of 2 April 1991 explains why the two adjoints to restriction, together with their exactness properties, should be treated as the essential structure \citep[pp.~8--9]{Grothendieck1991Letter}. The appendix examines the small categorical configurations behind that program. Its final gluing theorem belongs to the theory of homotopy colimits; its earlier counterexamples concern ordinary pushouts in \(\Cat\). Keeping the two operations distinct makes the apparent reversal between the examples and the closing conjecture intelligible.

There are also two different uses of asphericity for an inclusion. A functor \(U\to X\) is aspherical in the usual sense if every comma category \(U/x\) is weakly contractible. An open immersion in these pages is a full sieve, and the local condition used for it is imposed on \(U\to\overline U\), where \(\overline U\) is its cosieve closure. Objects outside \(\overline U\) have empty comma categories. This distinction will matter for the boundary of a subdivided square: that boundary has the homotopy type of a circle, even though the indicated local condition holds. We refer to the precise conventions of Chapter XII, \S V.2 \citep{Grothendieck1990XII}, and to the earlier companion's DER\_II-a, ms.~36, ``The quantifier in total asphericity,'' for the universal quantifier in that second notion.
\end{dcomment}

\begin{dcomment}{exposition}{p1:Le fait que $b$ soit un plus grand élément}{What the product actually contracts}
Grothendieck begins by asking how much of a terminal object can be replaced by the existence of maps. The objects $b$ and $c$ are \emph{weakly terminal}: every object admits a map to each of them, without any assertion of uniqueness. Their product $a=b\times c$ is weakly terminal as well. The argument establishes the exact conclusion needed at the first stage: the nerve $NX$ is connected and simply connected. Here and below $BX=|NX|$, and ``aspherical'' in the terminology of these pages means weakly contractible, not merely a $K(\pi,1)$.

The cancellation argument can be compressed without losing its content. For two maps $s,t:x\to a$, insert the mixed pair $(p_bs,p_ct)$. Cancellation against the now invertible projections shows that all three maps have the same image in the groupoid completion $\Pi X$. Choose one such image $\alpha_x:x\to a$ for each object, taking $\alpha_a=1_a$. For every arrow $u:x\to y$, both $\alpha_yu$ and $\alpha_x$ come from arrows $x\to a$ and therefore agree. The same identity then holds for inverse arrows and their composites. Every morphism $x\to a$ in $\Pi X$ is consequently $\alpha_x$. Thus $a$ is terminal in this groupoid, which is equivalent to the terminal groupoid. The longer zigzag induction in the text is a complete proof of this fact; its undefined $a_0$ and $u_0$ are printing slips, not a missing argument.

There is a significant distinction between this proposition and the later four-object example. In that example $b\amalg_a c=1$ in the presheaf category, but $a$ is \emph{not} $b\times c$: at the added object there are two maps to $a$ and only one pair of maps to $(b,c)$. It cannot by itself refute the opening product proposition. The finite two-object category below does refute that proposition, although its first obstruction occurs in degree three. We give the calculation because the provisional degree-two calculation in the manuscript does not survive verification.
\end{dcomment}

\begin{dcomment}{exposition}{p3:J'admet ici (provisoirement) ceci}{The finite-power category and the missing symmetries}
Grothendieck now makes the product condition explicit in a category small enough to calculate. Write $X_G$ for the category with objects $1,2$, representing $b,b^2$, and
\[
 \Hom_{X_G}(m,n)=\bigl(\{1,\ldots,m\}\times G\bigr)^n.
\]
A morphism lists, for each output coordinate, one input coordinate and one element of the monoid $G$; composition substitutes these expressions. Thus $(i,g):b^m\to b$ means $g\circ p_i$. This convention fixes multiplication even when $G$ is not commutative. The evident maps $2\to1$ exhibit $2$ as the product of $1$ with itself. Both objects map to $1$, so the hypotheses of the opening proposition hold. For a group $G$, the opposite category is the full category of free $G$-sets of ranks one and two.

Let $S$ be the full subcategory of finite sets on $\mathbf1,\mathbf2$, and put $T=S^{\mathrm{op}}$. The functor $X_G\to T$ forgets the monoid labels. A map $v:m\to n$ of $T$, equivalently a function $n\to m$, induces the coordinate-selection homomorphism $G^m\to G^n$. The arrows with all labels equal to $1_G$ are chosen cocartesian lifts, so $X_G$ is the Grothendieck construction of $n\mapsto BG^n$, where here $BG^n$ denotes the one-object category. This proves the asserted opfibration and justifies the cohomological spectral sequence, with constant abelian coefficients $M$:
\[
 E_2^{p,q}=R^p\!\varprojlim_{S} F_q
 \ \Longrightarrow\ H^{p+q}(BX_G;M),
 \qquad F_q(n)=H^q(BG^n;M).
\]
The functor $F_q$ is covariant on $S$. A modern treatment of this spectral sequence for Grothendieck constructions is \citet[Corollaries 3--4]{Cegarra2020}; the underlying homotopy-colimit comparison is \citet[Theorem 1.2]{Thomason1979}.

The provisional formula is already contradicted by a constant diagram $F=M$. The category $T$ has an initial object, hence contractible nerve, and $H^{>0}(T,M)=0$. Yet the proposed differential is $1_M-1_M=0$, whose cokernel is $M$. Changing the printed $H^2$ to the evidently intended $H^1$ does not repair this. The full category $S$ contains two inclusions, their common retraction, and the transposition of $\mathbf2$; retaining only a parallel pair loses relations and higher derived limits. On the other hand the formula for $H^0$ is correct: a compatible pair is determined by its value at $\mathbf1$, and the equality of the two induced values at $\mathbf2$ implies compatibility with the retraction and transposition.
\end{dcomment}

\begin{dcomment}{development}{p3:sont les transposés de}{A small resolution for the actual base category}
The point of repairing the base calculation is to retain Grothendieck's finite example, while locating its obstruction in the correct degree. The following elementary computation supplies all the derived limits used below. We did not find this calculation for the appendix's two-object category in the literature checked; the argument is included in full, without a claim of priority.

\textbf{Derived-limit calculation.}\ours
Let $F:S\to\mathrm{Ab}$, let $i_1,i_2:\mathbf1\to\mathbf2$ be the inclusions, let $r:\mathbf2\to\mathbf1$ be the unique map, and let $\tau$ transpose $\mathbf2$. Put
\[
 A=F(\mathbf1),\quad B=F(\mathbf2),\quad
 L=\ker(F(r):B\to A),\quad
 L^- =\{v\in L:F(\tau)v=-v\}.
\]
Then
\[
 \begin{aligned}
 R^0\!\varprojlim_S F&=\ker(F(i_1)-F(i_2):A\to B),\\
 R^1\!\varprojlim_S F&=L^-/(F(i_1)-F(i_2))A,\\
 R^p\!\varprojlim_S F&=\operatorname{Ext}^{p-1}_{\mathbb Z[C_2]}
       (\mathbb Z_{\mathrm{sgn}},L)\qquad(p\ge2).
 \end{aligned}
\]
The action of $C_2$ on $L$ is induced by $F(\tau)$.

\emph{Proof.} Work in the abelian category of covariant functors on $S$. The functors $P_j=\mathbb Z[S(\mathbf j,-)]$ are projective, since $\operatorname{Hom}(P_j,F)=F(\mathbf j)$ and evaluation is exact. The augmentation $P_1\to\underline{\mathbb Z}$ gives
\[
 0\longrightarrow K\longrightarrow P_1
   \longrightarrow\underline{\mathbb Z}\longrightarrow0.
\]
Here $K(\mathbf1)=0$, and $K(\mathbf2)=\mathbb Z(i_1-i_2)$ is the sign representation; the two constant endomorphisms act as zero. This immediately gives the formula for degree zero and identifies the degree-one group as the stated cokernel, because a map $K\to F$ is precisely an element of $L^-$.

For completeness the higher Ext groups can be computed inside the same tiny category. Precomposition with $r$ gives a monomorphism $P_1\to P_2$. It has the natural retraction induced by $i_1$, because $ri_1=1$. Consequently $Q=P_2/P_1$ is projective. It is zero at $\mathbf1$ and is the regular representation $\mathbb Z[C_2]$ at $\mathbf2$, with basis the classes of $1$ and $\tau$; both constant maps act as zero. Any free $\mathbb Z[C_2]$-resolution of the sign representation therefore gives a projective resolution of $K$, by replacing the regular module with $Q$. Applying $\operatorname{Hom}(-,F)$ gives exactly the complex computing $\operatorname{Ext}^*_{\mathbb Z[C_2]}(\mathbb Z_{\mathrm{sgn}},L)$: maps out of $Q$ are elements annihilated by $F(r)$. Finally dimension shifting through the displayed short exact sequence proves the formula for $p\ge2$. This also proves that no omitted higher relation is being assumed away.\hfill$\square$

For example, take $F(\mathbf1)=0$, $F(\mathbf2)=\mathbb Z$ with $\tau$ acting as $-1$. The formula gives $R^3\!\varprojlim F\cong\mathbb Z/2$, from the usual alternating resolution of $C_2$. Thus even the claimed vanishing in all degrees at least two is false for the full base category.
\end{dcomment}

\begin{dcomment}{development}{p4:Ainsi, si $G$ est un groupe abélien}{The counterexample survives, one degree later}
Grothendieck's proposed category is indeed nonaspherical for $G=C_2$. Its degree-two cohomology, however, is zero. The distinction matters twice: it changes the audit of the first four pages, and it supplies a sound finite example for the pushouts that follow.

\textbf{The actual low-degree calculation.}\ours
For every monoid $G$, the space $BX_G$ is $2$-connected. For $G=C_2$ one moreover has
\[
 H^3(BX_{C_2};\mathbb F_2)\cong\mathbb F_2.
\]
In particular the opening weak-contractibility assertion is false under its exact product hypotheses.

\emph{Proof.} The groupoid argument already proves simple connectivity. Use the spectral sequence and notation of the preceding two comments. The $q=0$ row is the constant diagram $M$, so its terms vanish for $p>0$. For $q>0$ the two coordinate pullbacks
\[
 H^q(BG;M)\rightrightarrows H^q(BG^2;M)
\]
have zero equalizer: restriction along $g\mapsto(g,1)$ sends the first pullback to the identity and the second to zero. Thus $E_2^{0,q}=0$ for $q>0$.

Set $A=H^1(BG;M)=\Hom_{\mon}(G,M)$, with $M$ regarded as an additive group. A normalized degree-one cocycle on the monoid is exactly an additive homomorphism, which proves this identification directly. A homomorphism $G^2\to M$ is uniquely the sum of its restrictions to the two factors, so the degree-one diagram is
\[
 F_1(\mathbf1)=A,\qquad F_1(\mathbf2)=A\oplus A,
\]
with $i_1,i_2$ inducing the two inclusions, $r$ inducing addition, and $\tau$ interchanging the factors. Hence $L=\{(a,-a):a\in A\}$, $L^-=L$, and the map $A\to L$ in the degree-one formula is an isomorphism. Thus $E_2^{1,1}=0$. All three entries of total degree two vanish; consequently $H^2(BX_G;M)=0$ for every abelian $M$. Simple connectivity and the universal coefficient theorem give $H^2(BX_G;M)=\Hom(H_2(BX_G;\mathbb Z),M)$. Taking $M=H_2(BX_G;\mathbb Z)$ forces this group to be zero. The Hurewicz theorem now gives $\pi_2(BX_G)=0$; see \citet[Theorems 3.2 and 4.32]{Hatcher2002}.

Now put $G=C_2$ and $M=\mathbb F_2=k$. The same degree-one diagram has $L=k$ with trivial $C_2$-action. The base calculation gives
\[
 E_2^{2,1}=\operatorname{Ext}^1_{\mathbb Z[C_2]}(\mathbb Z_{\mathrm{sgn}},k)=k.
\]
Indeed the free resolution of the sign module has differentials alternating between $1+\tau$ and $1-\tau$, both zero after applying $\operatorname{Hom}(-,k)$. No differential can kill this term: its sole possible incoming differential is from $E_2^{0,2}=0$ and its sole possible outgoing differential lands in $E_2^{4,0}=0$.

To identify the whole degree-three group, the only remaining possible entry is $E_2^{1,2}$. Write $H^*(BC_2;k)=k[t]$ and $H^*(B(C_2^2);k)=k[x,y]$, with all three generators in degree one; these are the usual projective-space and K\"unneth calculations \citep[Theorems 3.15 and 3.19]{Hatcher2002}. In degree two, $F_2(r)$ sends each of $x^2,xy,y^2$ to $t^2$. The $\tau$-fixed part of its kernel is therefore $k(x^2+y^2)$, exactly the image of $F_2(i_1)-F_2(i_2)$. The degree-one formula gives $E_2^{1,2}=0$. All other entries of total degree three vanish, proving the displayed isomorphism.\hfill$\square$

This proof also explains the discrepancy in the manuscript. The abstract cokernel of $a\mapsto(a,-a)$ is indeed $A$, but it is not the derived limit being computed. The transposition and the common retraction replace that cokernel by the preceding resolution. For $G=C_2$ the category has two objects and $2+4+4+16=26$ arrows. It supplies the finite counterexample promised in the text, with a degree-three obstruction. The stronger assertion in footnote~2 that the general product hypotheses need not imply $2$-connectivity requires a different example. The next comment verifies it with an actual product and a two-sphere.
\end{dcomment}

\begin{dcomment}{development}{p4:Ainsi, si $G$ est un groupe abélien}{A product whose category is a two-sphere}
The correction of the monoid calculation leaves one assertion in footnote~2 to check: the exact product hypotheses may fail to force $2$-connectivity. That assertion is true. One extra object supplies a smaller example in object count, and reveals geometrically the second homotopy group that the proposed monoid calculation did not produce. The construction below is our synthesis of the product condition and the elementary nerve-of-a-cone calculation; we did not locate this particular finite example in the literature searched, and make no claim of priority.

\textbf{A family realizing wedges of two-spheres.}\ours
For every integer $r\ge1$ there is a category $Y_r$ with three objects $d,b,a$, where $b$ is weakly terminal and $a=b\times b$, such that
\[
 BY_r\simeq\bigvee^{(r-1)(r-2)/2}S^2.
\]
The wedge with zero summands means a point. In particular $Y_3$ has $21$ arrows, satisfies the opening hypotheses with $c=b$, and has the homotopy type of $S^2$.

\emph{Proof.} Let $D$ be an $r$-element set. On the two objects $b=b^1$ and $a=b^2$ use the category $T=S^{\mathrm{op}}$ already encountered, with
\[
 \Hom_T(b^m,b^n)=\{1,\ldots,m\}^{\{1,\ldots,n\}},
 \qquad m,n\in\{1,2\}.
\]
These are coordinate-selection maps, composed by substitution. Add the object $d$, its identity, the maps
\[
 \Hom(d,b)=D,\qquad \Hom(d,a)=D^2,
\]
and no maps from $a$ or $b$ to $d$. Composition with old arrows again selects coordinates. Substitution is associative, so this defines a category. The product bijection
\[
 \Hom(z,a)\xrightarrow{\ (p_1,p_2)\ }
 \Hom(z,b)\times\Hom(z,b)
\]
holds for $z=a,b$ by the definition of $T$, and for $z=d$ it is the identity on $D^2$. Hence $a=b\times b$ in the enlarged category. As $D$ is nonempty, all three objects map to $b$.

To see the homotopy type, consider $E=d\backslash T$. An object of $E$ is a list of length one or two with entries in $D$; a morphism selects coordinates to obtain its target list. Let $P$ be the poset of nonempty subsets of $D$ of size at most two, ordered by reverse inclusion. Sending a list to its support defines $s:E\to P$. Choose one ordering of $D$ and let $j:P\to E$ list each subset without repetition in that order. Coordinate selection makes $j$ a functor and $sj=1_P$. There is a natural transformation $js\to1_E$: from the distinct list of a support to any given list, repeat the indicated coordinates. This map is unique because each label occurs once in its source. Its uniqueness also proves naturality for every coordinate-selection arrow. Therefore $BE\simeq BP$.

For clarity, the last gluing step can be read directly on the nerve. Every nondegenerate simplex of $NY_r$ is either wholly in $NT$, or starts at $d$ and continues in $T$. A simplex of the second kind is precisely the cone on a simplex of $N(d\backslash T)$. Its face opposite $d$ is obtained by forgetting the initial arrow from $d$. Thus $NY_r$ is the simplicial mapping cone of $NE\to NT$. After realization it is the mapping cone of $BE\to BT$. The category $T$ has initial object $b$, so $BT$ is contractible and
\[
 BY_r\simeq\Sigma BE\simeq\Sigma BP.
\]
The space $BP$ is the barycentric subdivision of the complete graph on the $r$ vertices of $D$. For $r\ge1$ that graph is connected and has $\binom r2-r+1=(r-1)(r-2)/2$ independent cycles: collapse a spanning tree to obtain the wedge of that many circles. Suspending gives the asserted wedge of two-spheres. Finally $T$ has $8$ arrows; adding $d$ contributes $1+r+r^2$ arrows. For $r=3$ the total is $21$.\hfill$\square$

This example verifies footnote~2 under the actual product assumptions, including its failure of $2$-connectivity. It also keeps the two repairs distinct. The two-object category $X_{C_2}$ first detects nonasphericity in degree three; the three-object category $Y_3$ detects it in degree two. The four-object category considered next in the manuscript has the latter homotopy type but satisfies only the terminal-presheaf-span condition.
\end{dcomment}

\begin{dcomment}{status}{p5:Je suspecte à présent que non}{Total asphericity without a terminal object}
The question concerns how far the useful gluing condition extends beyond representability. In the sense used here, a category \(A\) is totally \(\Winf\)-aspherical when, for every functor \(f:A\to C\), the presheaf \(f_!1_A\) has weakly contractible category of elements. This is the universal left-Kan-extension condition of Chapter XII, ms.~88--90; the editors explicitly distinguish it from the total asphericity of \emph{Pursuing Stacks}. The suspicion that a terminal object is necessary is false. Filtered categories provide counterexamples; indeed Grothendieck records this correction himself in Chapter XII, ms.~109, footnote~78 \citep{Grothendieck1990XII}. Chapter II, Corollaries~2.17--2.18 (ms.~34--35), gives the underlying filtered argument \citep{Grothendieck1990II}. We recall its mechanism because it answers the question on this page, while the earlier companion's DER\_II-a, ms.~35--36, already discusses its broader setting.

For any small \(A\), the canonical presentation of the terminal presheaf gives
\[
  f_!1_A\cong\mathop{\mathrm{colim}}_{a\in A}C(-,f(a)).
\]
Suppose that \(A\) is filtered, and denote this presheaf by \(F\). We verify directly that \(C/F\) is filtered. It is nonempty, since the identity of any \(f(a)\) supplies an object. Given \((c,x)\) and \((d,y)\), choose representatives \(u:c\to f(a)\) and \(v:d\to f(b)\). A common successor \(k\) of \(a,b\) produces arrows from both objects to \((f(k),[\id_{f(k)}])\). Finally, let \(r,s:(c,x)\rightrightarrows(d,y)\) be parallel arrows. Represent \(y\) by \(v:d\to f(b)\). The equality \(F(r)y=F(s)y\) says that \(vr\) and \(vs\) become equal at some later stage of the filtered colimit. Thus there is an arrow \(b\to k\) such that its image, composed with \(v\), equalizes \(r\) and \(s\). This is an equalizer after postcomposition inside \(C/F\).

The nerve of a filtered category is weakly contractible. One way to see this is to take the finitely many objects and arrows occurring in a simplicial representative of a homotopy class, choose a common target, and then equalize the finitely many competing composites. The resulting compatible cone gives a null homotopy. Equivalently, the usual finite-diagram characterization of filteredness cones off every finite simplicial diagram. Consequently \(C/F\) is weakly contractible for every \(f\), as required. The elementary eventual-equality fact used above is the standard filtered-colimit argument \citep[\S4.19]{StacksFiltered}.

For a concrete infinite example take \(A=(\mathbb N,\leq)\). It has no greatest element and hence no terminal object. The phenomenon is not a consequence of infinitude: the one-object category with endomorphism monoid \(\{1,e\}\), \(e^2=e\), is filtered, since postcomposition by \(e\) equalizes any pair of endomorphisms. Its unique object is not terminal, because it has two endomorphisms. For a functor out of this category, \(f_!1_A\) is the image of the induced idempotent on the represented presheaf \(C(-,f(*))\). Taking \(C=A\) and \(f=\id_A\) makes this image the terminal presheaf, with one element at the unique object, whereas the sole representable has two elements there. Thus the gluing weight need not be representable. These examples settle the question as asked; they do not assert an exhaustive classification of totally aspherical categories.
\end{dcomment}

\begin{dcomment}{examples}{p6:ne mettre en jeu que des catégories finies}{Adjoining a terminal object after the pushout}
This construction uses the failure of asphericity to test whether an ordinary pushout preserves weak equivalences. Its mechanism is transparent on representables. For any object $x$ of a category satisfying the opening hypotheses, put $B_x=\Hom(x,b)$ and $C_x=\Hom(x,c)$. Both sets are nonempty, and
\[
 B_x\amalg_{B_x\times C_x}C_x=\{*\}.
\]
Every element of $B_x$ is identified with every element of $C_x$, so there is exactly one equivalence class. As presheaf colimits are computed objectwise, the span $b\leftarrow b\times c\to c$ has terminal colimit. This verifies $p_!1_\Psi=1_{U'}$ without a homotopy assertion.

Write $\Psi^{\triangleright}$ for $\Psi$ with a new terminal object. In the pushout $U'\amalg_\Psi\Psi^{\triangleright}$, a map from $x\in U'$ to the new object is represented by a map from $x$ to an object of the span, modulo its two identifications. Thus the presheaf of maps to that object is precisely $p_!1_\Psi$, and the pushout is $(U')^{\triangleright}$. The original inclusion $\Psi\to\Psi^{\triangleright}$ is aspherical: the comma category over the new vertex is $\Psi$, with initial object, and the comma categories over old vertices have terminal objects. After the pushout, $U'\to(U')^{\triangleright}$ is a weak equivalence exactly when $BU'$ is weakly contractible.

Taking $U'=X_{C_2}$ therefore validates the announced failure. Both the original span and its cone are contractible, whereas $H^3(BU';\mathbb F_2)=\mathbb F_2$. The source's asserted obstruction in $\pi_2$ must be replaced by the degree-three calculation above; the categorical pushout and its terminal object require no alteration. This is also the point where the distinction from a homotopy pushout becomes visible: a homotopy pushout along the equivalence $\Psi\to\Psi^{\triangleright}$ would retain the homotopy type of $U'$, while the ordinary categorical pushout has become a cone.
\end{dcomment}

\begin{dcomment}{exposition}{p7:et même contractiles toutes trois}{Two cone vertices and a suspension}
Grothendieck strengthens the example by making the other map out of $\Psi$ a weak equivalence as well. Let $C=U'$, let $e$ be the terminal object adjoined to form $U''=C^{\triangleright}$, and retain $t$ for the cone vertex coming from $X=\Psi^{\triangleright}$. The left Kan extension $p'_!1_C$ is the presheaf with value one on the old objects of $C$ and value zero at $e$: no map runs from the new terminal object back to $C$. This is the subterminal presheaf called $U'$ in the displayed formula.

Consequently the pushout $X''$ consists of $C$ with two incomparable new terminal vertices $e,t$, and one map from each old object to each vertex. There are no arrows between $e$ and $t$. Every nondegenerate nerve simplex lies in one of the two cones; their intersection is exactly $NC$. Hence
\[
 BX''=C(BC)\cup_{BC}C(BC)\simeq\Sigma BC.
\]
This gives the omitted nonasphericity argument directly. For the finite choice $C=X_{C_2}$, suspension in reduced cohomology yields
\[
 \widetilde H^4(BX'';\mathbb F_2)
   \cong\widetilde H^3(BC;\mathbb F_2)
   \cong\mathbb F_2.
\]
Thus $X''$ is not weakly contractible, although $\Psi$, $X$ and $U''$ are contractible. Both maps to $X''$ from the latter two categories fail to be weak equivalences. No converse to Quillen's Theorem A is required. In general a noncontractible space may have contractible suspension if it is connected and acyclic, so a bare appeal to the nonasphericity of a comma category would not suffice. Here the explicit surviving cohomology class proves exactly what is needed.
\end{dcomment}

\begin{dcomment}{exposition}{p8:Je sais qu'il faut, pour un tel contre-exemple}{A presheaf pushout remembers only components}
Grothendieck now seeks a finite obstruction that can be inserted into a much larger category. The decisive distinction is between the ordinary pushout of three representable presheaves and the homotopy carried by its category of elements. We make that distinction explicit before using the example. Write $y(a)=\Hom_C(-,a)$, and let
\[
 L=y(b)\amalg_{y(a)}y(c)
\]
for a span $b\xleftarrow{\beta}a\xrightarrow{\gamma}c$ in $C$. For $x\in C$, form the bipartite multigraph $K_x$ with vertices $b(x)\amalg c(x)$ and one edge for every $u\in a(x)$, joining $\beta u$ to $\gamma u$. Then $L(x)=\pi_0K_x$. An ordinary pushout has forgotten the cycles in this graph.

This description proves the assertion about monomorphisms, including the case where $L$ is not terminal. Let $E$ have objects $(v,x,u)$, with $v\in\Psi$ and $u:x\to f(v)$ for the given span $f:\Psi\to C$. A morphism to $(v',x',u')$ consists of $h:v\to v'$ and $r:x\to x'$ satisfying $f(h)u=u'r$. The projection $E\to\Psi$ has right adjoint $v\mapsto(v,f(v),1)$, so $BE$ is contractible. There is a canonical functor
\[
 P:E\longrightarrow C/L,\qquad (v,x,u)\longmapsto(x,[u]).
\]
Here $C/L$ denotes the category of elements of the presheaf $L$. For $(x,l)\in C/L$, the comma category $(x,l)\backslash E$ retracts up to a natural transformation onto the incidence category of the component $l$ of $K_x$. Indeed an object of this comma category includes $r:x\to x'$; send it to $(v,x,u'r)$ and use the morphism given by $r$ for the comparison with the original object. The incidence category has one object for each vertex and edge, and two arrows from each edge to its endpoints.

If $\beta$ is monic, every vertex in $b(x)$ has degree at most one. Every connected component of $K_x$ is consequently a star, possibly a single vertex, and its incidence category is contractible. Quillen's Theorem A, in its comma-category form, gives $BP:BE\simeq B(C/L)$. Thus $L$ is aspherical. The same argument applies with $b$ and $c$ interchanged. This supplies the proof promised in the source; the only external ingredient is the comma-category criterion \citep[\S1, Theorem A]{Quillen1973}.

There is also a useful conclusion without a monomorphism hypothesis. If $L=e_C$, every $K_x$ is nonempty and connected, and $C$ is simply connected. To see this directly, work in $\Pi C$. Put $d=[\beta][\gamma]^{-1}:c\to b$. Send a vertex $v\in b(x)$ to $[v]:x\to b$ and a vertex $w\in c(x)$ to $d[w]$. The two endpoints of an edge $u\in a(x)$ have the same image, since $d[\gamma u]=[\beta u]$. Connectedness therefore gives a single arrow $s_x:x\to b$. For every $r:x\to y$, precomposition of a vertex of $K_y$ proves $s_y[r]=s_x$. These identities extend to all zigzags and force every arrow from $x$ to $y$ in $\Pi C$ to equal $s_y^{-1}s_x$. Thus $\Pi C$ is equivalent to the terminal groupoid. The correction on manuscript p.~10 is justified: the first possible obstruction is in degree two.
\end{dcomment}

\begin{dcomment}{examples}{p9:comme $U'$, pour les contre-exemples pages 6 et 7.}{The missing retraction and the two-sphere}
The four-object category is useful because its relations can be seen as attaching two disks to the same circle. We use the notation $Z$ of the following pages, suppressing subscripts: its objects are $0,a,b,c$, its generating arrows are $\lambda,\mu:0\rightrightarrows a$, $\beta:a\to b$, and $\gamma:a\to c$, and its relations identify the two composites to each of $b,c$. The earlier companion block ``The retracted proposition and a visible sphere'' (\texttt{DER\_II-a}, manuscript p.~36) already records the general obstruction to total asphericity. Here we need the precise attaching maps, because the next pushout will suspend this particular sphere once more.

Let $D=(0\rightrightarrows a)$. A retraction $Z\to D$ would fix $a$, so it must send $b$ to $a$ and $\beta$ to $1_a$. Applying it to $\beta\lambda=\beta\mu$ would give $\lambda=\mu$, a contradiction. This is the exact point at which the proposed retraction fails. The full subcategories on $\{0,a,b\}$ and $\{0,a,c\}$ are instead the two terminal cones on $D$. Every composable string in $Z$ belongs to one of them, and their intersection is $D$. Consequently
\[
 NZ=N(D^{\triangleright_b})\amalg_{ND}N(D^{\triangleright_c}),
 \qquad BZ\cong\Sigma BD\cong S^2.
\]
The first identity is an actual identity of simplicial sets. Realization preserves this union, and $BD$ is the circle with two vertices and two edges. This proves the correction already supplied by footnote~7, rather than replacing one guessed homotopy type by another.

The calculation is also visible in normalized chains. There are four vertices, six nonidentity arrows, four nondegenerate $2$-simplices, and no higher nondegenerate simplices. With $[u,v]$ denoting a composable pair, a fundamental cycle is
\[
 [\lambda,\beta]-[\mu,\beta]-[\lambda,\gamma]+[\mu,\gamma].
\]
Its boundary vanishes by the two defining relations; it generates $H_2(BZ;\mathbf Z)\cong\mathbf Z$. In particular $\pi_1(BZ)=0$ and $\pi_2(BZ)\cong\mathbf Z$.

The presheaf pushout nevertheless is terminal. Its pointwise calculation is
\[
\begin{array}{c|cccc}
 x&0&a&b&c\\\hline
 |a(x)|&2&1&0&0\\
 |b(x)|&1&1&1&0\\
 |c(x)|&1&1&0&1\\
 |L(x)|&1&1&1&1
\end{array}
\]
and all restriction maps of $L$ are the unique maps between singletons. Thus $Z/L\cong Z$, while $L=e_Z$. The same data establish that $\Psi$ is not totally $\Winf$-aspherical. Both arrows out of $a$ fail to be monic, and the same pair $\lambda,\mu$ witnesses both failures. The word ``universal'' refers to the free category with these specified generators and relations; arbitrary failures of the two monomorphisms need not have a common witnessing pair.
\end{dcomment}

\begin{dcomment}{exposition}{p11:qui a le type d'homotopie du cercle.)}{The link detects the failure of Dwyer's condition}
Grothendieck is locating the obstruction where it will survive a change of ambient category. For the closed inclusion $\Psi\hookrightarrow Z$, the relevant link is $0\backslash\Psi$. Its four objects are the two arrows $\lambda,\mu$ to $a$ and the unique arrows to $b,c$. Each of the first two objects maps to each of the last two, and there are no other nonidentity arrows. Its nerve is the square graph, hence its classifying space is $S^1$. Projection to $\Psi$ is the discrete opfibration classified by the covariant functor $\Hom_Z(0,-)|_\Psi$, with fibre cardinalities $2,1,1$.

This link has neither an initial object nor a contractible nerve. The inclusion is therefore neither a closed Dwyer inclusion nor coaspherical. These are separate conclusions: Dwyer's condition asks for a left adjoint to the inclusion into the smallest open neighbourhood, equivalently for initial objects in the relevant comma categories; weak contractibility alone would be a weaker requirement. For the complementary open inclusion $\{0\}\hookrightarrow Z$, the slice over $a$ is the discrete two-object category, so this inclusion also fails Dwyer's condition. The larger category can be simply connected while these local links retain nontrivial information. That local information, rather than the nonexistent retraction, will drive the finite pushout counterexample.
\end{dcomment}

\begin{dcomment}{exposition}{p12:et où $U$ est son ‘bord’}{The boundary and its categorical closure}
Grothendieck wants to know whether ordinary subdivisions of a square provide attaching maps compatible with weak equivalences. The notation carries two geometries: the realization of the nerve is a genuine square, while the open and closed subsets of the category are determined by the directions of its arrows. Write
\[
 \Phi_n=(0\longrightarrow1\longleftarrow2\longrightarrow\cdots
                \longrightarrow2n-1\longleftarrow2n),
 \qquad \dot\Phi_n=\{0,2n\}.
\]
This is the convention in Chapter XII, manuscript p.~114 \citep{Grothendieck1990XII}; the earlier companion's discussion of directed intervals (\texttt{DER\_VII-a}, manuscript pp.~1--19) explains its relation to the underlying unoriented interval. The boundary $U$ consists of the pairs with at least one coordinate an endpoint. It is a sieve, and
\[
 BX\cong[0,1]^2,\qquad BU\cong\partial[0,1]^2\cong S^1.
\]
In particular $U\to X$ is not a weak equivalence. The expression ``immersion ouverte $W$-asph\'erique'' in the proposition must be read in the local sense $U\to\overline U$: it is this induced functor to the categorical closure that is aspherical.

The slice calculation makes the distinction exact. Put
\[
 E_n=\{0,2n\},\qquad E_n^+=\{0,1,2n-1,2n\}.
\]
The cosieve generated by $U$ is
\[
 \overline U=(E_n^+\times\Phi_m)\cup(\Phi_n\times E_m^+).
\]
For $t\notin\overline U$, the slice $U/t$ is empty. For $t\in U$ it has terminal object $t$. Suppose that $t\in\overline U\setminus U$. If exactly one coordinate is adjacent to an endpoint and the other is an interior even vertex, then $U/t$ is a point. If that other coordinate is an odd vertex not adjacent to an endpoint, $U/t$ is the three-object interval with a terminal middle vertex. Finally, when both coordinates are adjacent to endpoints, $U/t$ is the five-object category $\widetilde\Psi$ drawn in the source. There are exactly four such $t$, because $n,m\geq2$.

The nerve of $\widetilde\Psi$ is the path with successive vertices $a',a,0,b,b'$. It is contractible, although neither $a$ nor $b$ receives a map from the other. Thus every slice over $\overline U$ is contractible, proving the claimed local asphericity, while the four exceptional slices have no terminal object. For a sieve of posets, Dwyer's condition is precisely the existence of these terminal objects over its generated cosieve: they assemble into the right adjoint to the inclusion. This proves both that $U\hookrightarrow X$ is not a Dwyer inclusion and that deleting the four exceptional vertices makes it one. Contractible slices suffice for a weak equivalence to the closure; an adjoint supplies the stronger structure that behaves well under arbitrary pushout.
\end{dcomment}

\begin{dcomment}{examples}{p15:dans le cas, disons, où $W \subseteq W_1$}{The pushout is a three-sphere}
The purpose of adjoining the new terminal object $e$ is to make the replacement $V'$ contractible without removing the obstruction in the link of the other cone point $t$. It is possible to calculate the entire pushout, so the argument need not depend on the converse excision principle mentioned in footnote~8. Rename $V'_0$ as $Z$ and the original three-object corner as $V$. Thus $Z$ is the four-object category just analysed, $V'=Z^{\triangleright_e}$, and $\overline V=V^{\triangleright_t}$.

Let $Q=V'\amalg_V\overline V$. Its objects are the four objects of $Z$ and two further objects $e,t$. It contains $Z$ as a full subcategory. For each $z\in Z$ there is exactly one arrow $z\to e$ and exactly one arrow $z\to t$; there are no arrows between $e$ and $t$. To check the only nontrivial assertion, the cocone over $V$ supplies the maps from its three objects to $t$. The two composites from the extra object $0'$ to $t$ agree, because the relation $\alpha\lambda=\alpha\mu$ already identifies their composites through $a$. Hence there is a unique arrow $0'\to t$. This description has the pushout universal property: a functor out of $Q$ is exactly a functor out of $Z$ together with one cocone over $Z$ with vertex the image of $e$ and one with vertex the image of $t$; the latter is uniquely determined by its restriction to $V$.

Every simplex of $NQ$ lies in one of the two cones on $NZ$, and their intersection is $NZ$. Therefore
\[
 BQ\cong\Sigma BZ\cong S^3,
 \qquad \widetilde H_k(BQ;\mathbf Z)=
 \begin{cases}\mathbf Z,&k=3,\\0,&k\ne3.\end{cases}
\]
Since $BV$, $B\overline V$, and $BV'$ are contractible, the maps $i$ and $p$ are weak equivalences, but $i'$ and $\overline p$ are not. This is a completely finite example: $Q$ has six objects, fourteen nonidentity arrows, sixteen nondegenerate $2$-simplices, and eight nondegenerate $3$-simplices. Its nerve has no higher nondegenerate simplices.

The sentence $\pi_1(V'_0)\cong\mathbf Z$ is present in the edition and must remain in the historical text. Its correct replacement in the argument is
\[
 \pi_1(V'_0)=0,\qquad \pi_2(V'_0)\cong\mathbf Z,
 \qquad \pi_3(Q)\cong\mathbf Z.
\]
The example proves the required assertion for $W=\Winf$. It does not prove the printed extension to every $W\subseteq W_1$: both $V'\to Q$ and $\overline V\to Q$ are equivalences on $2$-types. Here $W_k$ means equivalence on homotopy groups in degrees at most $k$ and on components.

There is a second, logically independent issue. For an arbitrary open inclusion with one-point complement, membership in $\Winf$ does not force its comma category over that point to be contractible. The referenced Chapter XII corollary explicitly assumes the inverse part of $W(7\,\mathrm{bis})$. The acyclic-link construction in the commentary on manuscript pp.~17--18 exhibits the failure of that inverse implication for spaces. In the present example both the link and the pushout are known spheres, so the conclusion follows from the displayed nerve calculation with no appeal to this extra axiom.
\end{dcomment}

\begin{dcomment}{exposition}{p16:ce qui montre bien que $V \longrightarrow X$}{Carrying the local obstruction to the boundary}
The last step places the six-object defect inside the subdivided rectangle. Grothendieck has chosen the corner so that the two surrounding inclusions are closed Dwyer inclusions. We record their adjoints, because they explain exactly which pushouts preserve homotopy and complete the logical step in the proof.

Use coordinates $0=(0,0)$, $a=(1,0)$, $b=(0,1)$, and $t_1=(1,1)$. The smallest open neighbourhood of $V=\{0,a,b\}$ in $U$ consists of these three points and $a'=(2,0), b'=(0,2)$. Its inclusion of $V$ has the left adjoint fixing $V$ and sending $a'$ to $a$ and $b'$ to $b$. Indeed these images are the least points of $V$ above the corresponding objects. The inclusion $V\hookrightarrow U$ is a cosieve, so it is a closed Dwyer inclusion.

Similarly $\overline V=\{0,1\}\times\{0,1\}$ is a cosieve of $X$. Its smallest open neighbourhood is $\{0,1,2\}^2$. Define $r_0(0)=0$, $r_0(1)=r_0(2)=1$; then $r_0\times r_0$ is left adjoint to the inclusion of $\overline V$ into this neighbourhood. Thus $\overline V\hookrightarrow X$ is also a closed Dwyer inclusion. This gives the source's five comma categories explicitly: $a'\backslash\overline V$ and $b'\backslash\overline V$ are two-object chains, whereas the other three are points.

Thomason's pushout theorem, applied after taking opposite categories, now identifies the nerves of the relevant categorical pushouts with homotopy pushouts \citep[Proposition~4.3 and Corollary~4.4]{Thomason1980}. One additional observation is needed for the reverse implication written in the source: $\overline V\to X$ is itself a weak equivalence, since both classifying spaces are contractible. Consequently its pushout $Q\to X'$ is a weak equivalence. In the commutative square
\[
\begin{tikzcd}
\overline V\arrow[r,"\sim"]\arrow[d,"\overline p"']&X\arrow[d,"q"]\\
Q\arrow[r,"\sim"]&X'
\end{tikzcd}
\]
two-out-of-three gives $\overline p\in\Winf$ if and only if $q\in\Winf$. Since $BQ\simeq S^3$, we obtain $BX'\simeq S^3$ and $q\notin\Winf$. A homotopy pushout by itself would not justify the converse: the two horizontal equivalences are the missing reason in this particular square.

Finally let $U'=U\amalg_VV'$. The other closed Dwyer inclusion, $V\to U$, ensures that $q_U:U\to U'$ is a weak equivalence because $p:V\to V'$ is one. Associativity of ordinary pushout identifies $X\amalg_UU'$ with the same $X'$. Thus the cobase-change functor along $U\to X$ takes the weak equivalence $q_U$ to the map $q$ which is not a weak equivalence. This proves exactly the definition of failure of a $W$-cofibration used in Chapter XII: the functor $U\backslash\Cat\to X\backslash\Cat$ fails to preserve weak equivalences. In modern terminology, this inclusion is not an $h$-cofibration. The target replacement still satisfies $BU'\simeq S^1$; the change in the ambient category is detected by its new third homology group.
\end{dcomment}

\begin{dcomment}{later}{p16:Cela rend explicite que $U \longrightarrow X$}{Dwyer maps, homotopy pushouts, and products}
Grothendieck's distinction between contractible links and adjoint neighbourhoods remains useful in the modern theory. Thomason's original Dwyer condition guarantees that the comparison from a pushout of nerves to the nerve of a categorical pushout is a weak homotopy equivalence. It therefore controls the entire change-of-base operation, rather than merely the homotopy type of the initial inclusion. One should not identify Dwyer maps with all cofibrations in the Thomason model structure: \citet[Lemma~4 and Proposition~2]{Cisinski1999Dwyer} corrected the claim that the original Dwyer class is closed under retracts. His larger class of pseudo-Dwyer maps retains the required pushout property.

The later result of \citet[Theorems~1.5--1.6]{HackneyOzornovaRiehlRovelli2024} strengthens the comparison to an equivalence of $(\infty,1)$-categories. Thus pushouts along Dwyer maps preserve more than the classifying-space homotopy type. Our two uses of closed Dwyer maps are obtained by passing to opposites; the reversal changes the direction of the neighbourhood adjunction but does not change the classifying-space conclusion.

The negative product phenomenon also has a precise later form. \citet[Definition~2.6]{Pavlov2024Thomason} uses $h$-cofibration for the preservation condition Grothendieck calls $W$-cofibration. In the same paper, Proposition~4.3 proves that no model structure on $\Cat$ with Thomason weak equivalences can be cartesian. Its proof uses a different finite cube construction, credited in Proposition~4.2 to Maltsiniotis. The present appendix supplies a particularly transparent earlier obstruction: the boundary of a product of alternating intervals can have contractible links in its closure and still fail the preservation condition. These are results about strict categorical pushouts; homotopy pushouts of spaces retain their usual invariance.
\end{dcomment}

\begin{dcomment}{development}{p16:Cela rend explicite que $U \longrightarrow X$}{The same defect on higher cubical boundaries}
The rectangle argument is local enough to survive in every higher dimension. This gives a concrete continuation of Grothendieck's construction: the defect need not grow with the dimension of the ambient cube. The literature checked for this extension comprised Thomason's pushout theorem, Cisinski's correction, the Dwyer pushout comparison of Hackney--Ozornova--Riehl--Rovelli, and Pavlov's distinct cube obstruction. We did not locate the following precise family there; the proposition is an elementary extension of the appendix, without a claim of priority.

\textbf{Proposition (one fixed obstruction in every dimension).}\ours
Let $d\geq2$, let $n_j\geq2$, and put
\[
 X_d=\prod_{j=1}^d\Phi_{n_j},\qquad
 U_d=\bigcup_{j=1}^d
   \Phi_{n_1}\times\cdots\times\dot\Phi_{n_j}\times\cdots\times\Phi_{n_d}.
\]
There is a finite weak equivalence $U_d\to U'_d$ whose pushout along $U_d\hookrightarrow X_d$ has target $X'_d$ with $BX'_d\simeq S^3$. In particular every inclusion $U_d\hookrightarrow X_d$ fails to be a $\Winf$-cofibration.

\emph{Proof.} Fix the sink $1$ in each of the last $d-2$ factors. Let $V$ be the three-object full subposet with first two coordinates $(0,0),(1,0),(0,1)$ and all other coordinates $1$. Let $S$ add the point $(1,1,1,\ldots,1)$. Thus $V\cong\Psi$ and $S\cong[1]\times[1]$. Every point of $V$ lies in $U_d$, while the fourth point of $S$ does not. Since each coordinate $1$ is maximal in its alternating interval, all successors of $S$ remain in $S$; all successors of $V$ that remain in $U_d$ remain in $V$. Hence $S\hookrightarrow X_d$ and $V\hookrightarrow U_d$ are closed inclusions.

Their Dwyer structures can be given without induction. The smallest open neighbourhood of $S$ is $\{0,1,2\}^d$. Map its first two coordinates by $r_0(0)=0$, $r_0(1)=r_0(2)=1$, and map every remaining coordinate to $1$. This sends an object to the least element of $S$ above it, and is therefore left adjoint to the inclusion of $S$. The smallest open neighbourhood of $V$ in $U_d$ is
\[
 \bigl(\{0,1,2\}\times\{0\}\;\cup\;
       \{0\}\times\{0,1,2\}\bigr)\times\{0,1,2\}^{d-2}.
\]
The same map lands in $V$ and is its left adjoint. Both inclusions are closed Dwyer inclusions.

Use the finite weak equivalence $p:V\to V'=Z^{\triangleright_e}$ constructed above. Put
\[
 U'_d=U_d\amalg_VV',\quad Q=S\amalg_VV',\quad
 X'_d=X_d\amalg_VV'.
\]
The first Dwyer pushout gives $U_d\simeq U'_d$. Both $BS$ and $BX_d$ are contractible, so $S\to X_d$ is a weak equivalence as well as a closed Dwyer inclusion. Its pushout gives a weak equivalence $Q\to X'_d$. The calculation on manuscript p.~15 identifies $BQ\simeq S^3$, so $BX'_d\simeq S^3$. Finally associativity of pushouts gives $X'_d=X_d\amalg_{U_d}U'_d$. The resulting map from the contractible $BX_d$ to $BX'_d$ is not a weak equivalence, proving the assertion. All categories involved are finite, since the displayed pushouts are obtained by the same finite gluing constructions.\hfill$\square$
\end{dcomment}

\begin{dcomment}{exposition}{p17:sous-catégorie pleine des flèches de source}{The arrows which join the two pieces}
Grothendieck now turns from unsuccessful pushouts to a decomposition which retains the attaching data. An open subcategory $A=X_0$ is a full sieve; its complementary closed subcategory $B=X_1$ is a full cosieve. Thus every object belongs to exactly one piece, and there are no arrows from $B$ to $A$. The additional information in $X$ is the functor
\[
 M:A^{\mathrm{op}}\times B\longrightarrow\Ens,
 \qquad M(a,b)=\Hom_X(a,b).
\]
In the language of distributors, $X$ is the collage of $M$; this description, and its simplicial counterpart, are given in \citet[\S\S7.1--7.2, pp.~309--311]{Joyal2008Quasi}. The direction of these attaching arrows is the feature that permits the argument below.

Write $E=\widetilde X_{01}$. Its objects are arrows $f:a\to b$ with $a\in A$ and $b\in B$; a morphism $f\to f'$ is a commuting square
\[
 (u,v):f\longrightarrow f',\qquad
 u:a\to a',\quad v:b\to b',\quad vf=f'u.
\]
Both vertical sides point in their ordinary directions. In particular $E$ is the ordinary crossing-arrow category, not a twisted-arrow category. The projection $p:E\to A$ is a split Grothendieck fibration: pullback along $u:a\to a'$ sends $f':a'\to b$ to $f'u:a\to b$. Its fibre is $a\backslash B$. Dually, $q:E\to B$ is a split opfibration, with fibre $A/b$ and transport given by postcomposition. The arrows $f$ themselves constitute a natural transformation $ip\Rightarrow jq$.

After realization this natural transformation provides a specified homotopy between the two maps $BE\rightrightarrows BX$, where $BC=|NC|$. Consequently it defines a map from the double mapping cylinder of $BA\leftarrow BE\to BB$ to $BX$. This is the comparison whose weak equivalence is at issue. The ordinary square does not commute strictly, and replacing its specified homotopy by an arbitrary one would change the question.

The easy cases stated in the manuscript have a useful explanation. If $i$ is aspherical, every $A/b$ is contractible, hence $q$ is a weak equivalence; the homotopy pushout then has the homotopy type of $BA$, which maps equivalently to $BX$. If $j$ is coaspherical, the dual argument uses the fibres $a\backslash B$ of $p$. The general argument must retain these fibres rather than assume their contractibility.
\end{dcomment}

\begin{dcomment}{development}{p17}{Acyclic links and a false converse}
The proposed $\mathbf W_\infty$ gluing formula and the opposite-side reflection invoked by the manuscript's ``i.e.'' have different strengths. The formula is true. Footnote~9 adds Loc~$(6)$, a numbered hypothesis which requires qualification. If it uses the full reversible convention of Chapter~XVI, Corollary~3.7, then cocartesianness together with $i\in W$ does imply $q\in W$, and dually for $j,p$: this is the additional condition Loc~$(6\mathrm C)$ \citep[\S3, Corollary~3.7, ms.~26]{GrothendieckDerivateursXVI}. Under that interpretation the conditional inference is valid, but its specialization to $\mathbf W_\infty$ is incompatible with the axiom. The following construction proves the failure for ordinary weak homotopy equivalences.

\textbf{Proposition (a contractible union with a noncontractible link).}\ours
There is a finite poset $X$, a one-object sieve $A=\{a\}$, and its complementary cosieve $B$, such that $BA$, $BB$, and $BX$ are contractible, whereas $B\widetilde X_{01}$ is not contractible. In particular neither $p$ nor $q$ need be a weak equivalence, even when both $i$ and $j$ are weak equivalences. Moreover, $\mathbf W_\infty$ fails Loc~$(6\mathrm b)$ in the convention of Chapter~XVI.

\emph{Proof.} We first choose a connected acyclic finite complex $P$ with nontrivial fundamental group. A standard source of such complexes is \citet[Example~2.38]{Hatcher2002}. For explicitness, take the presentation complex
\[
 \langle s,t\mid s^3(st)^{-2},\ t^5(st)^{-2}\rangle.
\]
It has one vertex, two edges and two two-cells. Its cellular differential in positive degrees is
\[
 0\longrightarrow\mathbb Z^2
 \xrightarrow{\left(\begin{smallmatrix}1&-2\\-2&3\end{smallmatrix}\right)}
 \mathbb Z^2\longrightarrow0.
\]
The determinant is $-1$, so its reduced integral homology vanishes. The assignments $s\mapsto(1\,2\,3)$ and $t\mapsto(1\,4\,3\,2\,5)$ define a nontrivial homomorphism to $A_5$: their product is $(1\,4)(2\,5)$, and both relations hold. Thus $P$ is not simply connected. By the finite form of the simplicial replacement theorem \citep[Theorem~2C.5]{Hatcher2002}, choose a finite simplicial complex $L\simeq P$ and let $K$ be its poset of nonempty faces. Its order complex is the barycentric subdivision of $L$, so $BK\simeq P$.

Adjoin two incomparable new elements $a,b$ to $K$, each smaller than every element of $K$, and call the resulting poset $X$. Set $A=\{a\}$ and $B=\{b\}\cup K$. Then $A$ is a sieve and $B$ is its complementary cosieve; $b$ is initial in $B$. The nerve of $X$ is the union of two cones on $NK$, along $NK$ itself. Hence
\[
 BX\cong\Sigma BK.
\]
This suspension is simply connected because $BK$ is connected, and its reduced homology vanishes by the suspension isomorphism. The simply connected homology form of Whitehead's theorem therefore makes it contractible \citep[Corollary~4.33]{Hatcher2002}.

The crossing arrows are exactly the unique arrows $a\to k$, for $k\in K$. Their commuting squares are exactly the order relations of $K$. Thus $E\cong K$, $p:K\to\{a\}$, and $q:K\hookrightarrow\{b\}\cup K$. Neither map is a weak equivalence, since $\pi_1(BK)$ is nontrivial.

For the last assertion, consider the full sieves
\[
 U=(\{a\}\cup K)^{\mathrm{op}},\qquad
 V=(\{b\}\cup K)^{\mathrm{op}}
 \quad\text{in }X^{\mathrm{op}}.
\]
They cover $X^{\mathrm{op}}$ and intersect in $K^{\mathrm{op}}$. Map their Mayer--Vietoris square to the square whose four vertices are the terminal category. The component maps from $U,V,X^{\mathrm{op}}$ are weak equivalences: their classifying spaces are two cones and the contractible suspension. The component from $K^{\mathrm{op}}$ is not. This violates precisely the inverse cube implication Loc~$(6\mathrm b)$; the full axiom includes that implication \citep[\S3, ms.~$16'$--$17'$, footnote~13]{GrothendieckDerivateursXVI}.\hfill$\square$

For $\mathbf W_\infty$, this also answers negatively the subsidiary question with $X_0\simeq e$ and $X,X_1$ aspherical. Its intended comma category is $a\backslash X_1=K$; the printed $a\backslash X$ would contain the initial object $\id_a$ and would be automatically contractible. The example does not contradict the gluing conjecture: its homotopy pushout is exactly the contractible suspension $\Sigma BK$. It disproves opposite-side reflection for $\mathbf W_\infty$, while leaving the stronger conditional principle intact.

The axiom labels cannot be identified solely by their numbers. Chapter~XV, \S9, ms.~117, calls a fibration axiom Loc~$(6)$ \citep{Grothendieck1990XV} and refers to Chapter~XII, \S IV, ms.~53--59; there, variant $\mathrm a'$ is a pullback condition for $W_\omega$-perfect functors \citep{Grothendieck1990XIIa}. Chapter~XVI instead writes the forward Mayer--Vietoris condition as Loc~$(6'\mathrm a)$, whereas the appendix's last line prints Loc~$(6\mathrm a')$ \citep[\S3, ms.~$17'$]{GrothendieckDerivateursXVI}. We retain these distinct labels and cannot certify that footnote~9 invokes the later reversible convention. We checked the standard acyclic-complex construction and the homotopy-colimit sources cited below; we did not locate this precise finite-poset application there. The marker records our worked application, without a claim of priority.
\end{dcomment}

\begin{dcomment}{development}{p18}{The open--closed gluing formula}
The final conjecture $\varphi\in\mathbf W_\infty$ has an affirmative answer. Its proof comes from the fact that a simplex in $NX$ crosses from $A$ to $B$ at most once. We give the precise specialization of Thomason's homotopy-colimit comparison that settles this question. We did not locate this exact formulation in the checked sources; the result below is presented as a consequence of an established theorem, without a claim of priority.

\textbf{Proposition (directed gluing).}\ours
For every small category $X$ with a full sieve $A$ and complementary full cosieve $B$, let $E$ be its crossing-arrow category. The natural transformation $ip\Rightarrow jq$ induces a weak homotopy equivalence
\[
 \operatorname{hocolim}(BA\longleftarrow BE\longrightarrow BB)
 \longrightarrow BX.
\]
Consequently the comparison $\varphi:\mathcal X\to X$ of the manuscript belongs to $\mathbf W_\infty$.

\emph{Proof.} To keep track of the unique crossing, introduce the bisimplicial set
\[
 S_{r,s}=\left\{
 a_0\longrightarrow\cdots\longrightarrow a_r
 \xrightarrow{f}b_0\longrightarrow\cdots\longrightarrow b_s
 \ \middle|\ a_i\in A,\ b_j\in B
 \right\},\qquad r,s\geq0.
\]
Its faces compose or omit arrows on the indicated side, and its degeneracies insert identities. The last horizontal face precomposes $f$; the first vertical face postcomposes $f$. Thus $S$ is the simplicial replacement of the contravariant functor
\[
 A^{\mathrm{op}}\longrightarrow\Cat,\qquad
 a\longmapsto a\backslash B.
\]
Its Grothendieck construction, with cartesian variance, is exactly $E$: a morphism over $u:a\to a'$ is a map $f\to f'u$ in $a\backslash B$, namely a $v$ with $vf=f'u$.

There is a comparison $\theta:\operatorname{diag}S\to NE$. On an $n$-simplex it gives the chain of crossing arrows
\[
 a_i\longrightarrow a_n\xrightarrow{f}b_0
 \longrightarrow b_i,\qquad 0\leq i\leq n,
\]
whose successive morphisms have sides $a_i\to a_{i+1}$ and $b_i\to b_{i+1}$. This is Thomason's comparison for the preceding contravariant functor. It is a weak equivalence; the precise variance and this formula are verified in \citet[\S3.1, equation~(9), and Theorem~3]{Cegarra2020}, which gives an independent proof of the comparison originally due to \citet[Theorem~1.2]{Thomason1979}. Both endpoint maps commute with $\theta$: they simply forget the $b$-chain or the $a$-chain.

It remains to identify what the mixed simplices assemble to. Let $R=\lVert S\rVert$ be the double geometric realization. It has maps $r_A:R\to BA$ and $r_B:R\to BB$ obtained by forgetting one side. The natural diagonal-to-double-realization homeomorphism identifies $R$ with $|\operatorname{diag}S|$. Form
\[
 D(R)=\bigl(BA\amalg(R\times[0,1])\amalg BB\bigr)
       /\bigl((z,0)\sim r_A(z),\ (z,1)\sim r_B(z)\bigr).
\]
For a mixed simplex as above, send $(u,v,t)\in\Delta^r\times\Delta^s\times[0,1]$ to the point of $\Delta^{r+s+1}$ with coordinates
\[
 ((1-t)u_0,\ldots,(1-t)u_r,tv_0,\ldots,tv_s).
\]
This respects faces and degeneracies, and respects the endpoint identifications. It supplies a homeomorphism $D(R)\cong BX$. Indeed, in a point with positive total weight on both pieces, $t$ is the total weight on the $B$ vertices; dividing the two sets of coordinates by $1-t$ and $t$ recovers $u,v$. At $t=0,1$ precisely the identifications defining $D(R)$ occur. This description is compatible with the realization quotient topology because it is the usual simplex-join homeomorphism on every simplex. Pure simplices give the two endpoint spaces.

The double mapping cylinder computes a homotopy pushout. Replacing its middle space $R$ by the weakly equivalent CW complex $BE$ therefore preserves its homotopy type, by homotopy invariance of this construction. The resulting map to $BX$ is the canonical one: the homotopy obtained from $ip\Rightarrow jq$, after composition with $\theta$, is homotopic relative to its endpoints to the displayed barycentric homotopy. On each prism, both lie in the same mixed simplex, and linear interpolation supplies compatible homotopies on its faces. Thus the canonical comparison, and not merely an abstract map between the same homotopy types, is a weak equivalence.

Finally, let $J$ be the span category with arrows $*\to0$ and $*\to1$, and take its covariant diagram with values $E,A,B$ and structure maps $p,q$. Its Grothendieck construction models $\mathcal X$. A comparison to $X$ sends its $A$- and $B$-objects to themselves and its $E$-object $f:a\to b$ to $a$; the arrows towards the $B$-part use $f$. Thomason's theorem identifies its classifying space with the preceding homotopy pushout, compatibly with this cocone. Hence $B\varphi$ is a weak equivalence.\hfill$\square$

The cases $A=\varnothing$ or $B=\varnothing$ are included: then $E$ is empty and the assertion reduces to the other piece. For every localizer $W$ containing $\mathbf W_\infty$, the comparison just proved also lies in $W$. If the span Grothendieck construction represents the $W$-homotopy pushout, this proves $W$-cocartesianness as well. For a fundamental localizer, Chapter~XVI explicitly supplies that interpretation under its forward axiom Loc~$(6'\mathrm a)$ \citep[\S3, ms.~24, footnote~15]{GrothendieckDerivateursXVI}. This comparison hypothesis requires no opposite-side reflection. The manuscript cites Loc~$(6\mathrm a')$ for its general-$W$ interpretation; because of the numbering ambiguity just documented, we do not silently identify that label with the distinct forward or full axioms of Chapter~XVI. The theorem for $\mathbf W_\infty$ is independent of that identification.

The proof also clarifies the preceding observation about a phenomenon peculiar to categories. Its realization is an explicit topological double mapping cylinder, and the mixed-simplex construction has a directed simplicial formulation. What a bare pair of complementary topological subsets does not supply is the specified category of attaching arrows.
\end{dcomment}

\begin{dcomment}{examples}{p18}{Why replacing arrows by all paths changes the answer}
Grothendieck's caution about paths concerns the attaching information, rather than an insufficiency of path spaces as models of homotopy types. A bridge of directed arrows records how the two pieces were assembled. A bridge of all paths in the already assembled space usually records a homotopy pullback instead. These two bridges can differ even in the smallest example with a loop.

Let $X$ have objects $a,b$ and precisely two nonidentity arrows $\lambda,\mu:a\rightrightarrows b$. Take $A=\{a\}$ and $B=\{b\}$. There are no nonidentity commuting squares between its two crossing arrows, so $E$ is the discrete two-point category. The nerve of $X$ has two vertices and two edges and no nondegenerate higher simplices. Thus
\[
 BX\cong S^1,\qquad
 \operatorname{hocolim}(*\longleftarrow\{\lambda,\mu\}\longrightarrow*)
 \cong S^1.
\]

Now replace the bridge by a model of the space of all paths from $a$ to $b$ in $BX$. Choosing one reference path identifies this space up to homotopy with $\Omega S^1$, whose components are contractible and indexed by $\mathbb Z$. The new homotopy pushout is therefore the graph with two vertices and one edge $e_k$ for each $k\in\mathbb Z$. Collapsing $e_0$ gives
\[
 \bigvee_{k\in\mathbb Z\setminus\{0\}}S^1.
\]
The comparison to $S^1$ sends the loop represented by $e_ke_0^{-1}$ to the integer $k$. Its map on fundamental groups is the homomorphism from the free group on these generators to $\mathbb Z$ which sends the generator indexed by $k$ to $k$. It is surjective but not injective; for example the word consisting of the generator indexed by $2$ followed by the inverse square of the generator indexed by $1$ lies in its kernel and is a nonempty reduced word. Thus even the $\mathbf W_1$ conclusion fails.

More generally, when a path category models $BA\times^h_{BX}BB$, the exception stated in the manuscript is valid: if $BA\to BX$ is a weak equivalence, its projection to $BB$ is a weak equivalence, so the associated homotopy pushout maps equivalently to $BX$; the other case is dual. This is a sufficient condition, not a characterization of every exceptional case. The excerpt specifies the objects of $\widetilde X'_{01}$ but does not spell out its morphisms. Our computation concerns the usual homotopically correct all-paths interpretation; it does not assert an equivalence for an unspecified raw category of paths.
\end{dcomment}

\begin{dcomment}{later}{p18}{From covering descent to higher local systems}
The last paragraph identifies the right source of confidence in the conjecture: a local system can be reconstructed from its restrictions to the two pieces and its transport across every attaching arrow. We can make the argument for coverings completely explicit before passing to higher coefficients.

Use covariant transport to identify $\operatorname{Rev}(C)$ with the category of functors $C\to\Ens$ which send every arrow to a bijection. An object of
\[
 \operatorname{Rev}(A)
 \mathbin{\mathop{\times}\limits^{2}_{\operatorname{Rev}(E)}}
 \operatorname{Rev}(B)
\]
is a pair $L_A,L_B$ with an isomorphism $\alpha:p^*L_A\xrightarrow{\sim}q^*L_B$. Define $L$ on the two full subcategories by $L_A,L_B$, and on a crossing arrow $f:a\to b$ by $L(f)=\alpha_f$. To check functoriality, there are only two mixed cases. The squares $(u,\id):fu\to f$ and $(\id,v):f\to vf$ in $E$ give
\[
 \alpha_{fu}=\alpha_f L_A(u),\qquad
 \alpha_{vf}=L_B(v)\alpha_f.
\]
They are exactly the two composition laws needed in $X$, since a composable chain cannot cross from $A$ to $B$ twice. All transports are bijective. A compatible pair of morphisms between such descent data gives a natural transformation of the reconstructed functors, and restriction is inverse to this construction. This proves the equivalence of covering categories in the manuscript; applying the same gluing description to the categorical mapping cylinder gives its displayed formula for $\operatorname{Rev}(\mathcal X)$.

This argument also identifies the fundamental groupoid directly. Its generators are the arrows internal to the two pieces and the crossing arrows, now formally invertible; the equations just verified supply all mixed composition relations. Consequently the comparison $\varphi$ induces an equivalence of fundamental groupoids. The source offers a sketch of this $\mathbf W_1$ assertion and reports a further sketch for $\mathbf W_2$; the preceding mixed-simplex proof establishes all degrees at once.

The contemporary interpretation of the proposed higher-gerbe argument uses local systems valued in spaces or in $n$-groupoids. If $\mathcal S$ denotes the $\infty$-category of spaces, define $\operatorname{Loc}(C)=\operatorname{Fun}(BC,\mathcal S)$, viewing $BC$ as its homotopy type. The proved pushout formula, by the universal property of a colimit, yields
\[
 \operatorname{Loc}(X)\simeq
 \operatorname{Loc}(A)\times^h_{\operatorname{Loc}(E)}
 \operatorname{Loc}(B).
\]
One can restrict the values to $n$-truncated spaces to obtain every finite stage. Straightening gives the precise realization of space-valued coefficients, and higher descent formalizes such pullbacks of coefficient categories; see \citet[Theorem~2.2.1.2, Corollary~3.3.4.6, and Theorem~6.1.3.9]{Lurie2009HTT}. Giraud's theory supplied the stack and gerbe language for the first nonabelian stages \citep{Giraud1971}; the claim here is about locally constant higher coefficients, not an identification of all stacks with local systems. The categorical question of the appendix is therefore settled without requiring a separate construction and comparison of $n$-gerbes for each $n$.
\end{dcomment}
